\section{Introduction}
Frank--Wolfe (FW) replaces projections with linear minimization oracles~\cite{braun}.
Backtracking makes its quadratic step model usable without a prescribed
curvature constant, but a small primal error need not be visible to an
observable stopping certificate. This paper asks when FW can certify
progress, and when a certificate can lead to exact termination.

We first study ordinary FW on a polytope whose entire minimizing set is
a full exposed face. Assume that $f$ is convex and continuously
differentiable on a neighborhood of $Q$, its gradient is Lipschitz on
all of $Q$, and the atom oracle is exact. Theorem~\ref{face:ratio}
proves that along every infinite nonterminal trajectory, the FW gap
divided by primal suboptimality tends to one. A normalized-atom-weight
argument bounds the total number of exceptional updates and yields the
local observable-certificate complexity in Corollary~\ref{face:local}.
This is a trajectory statement; it does not assert the same ratio for
every near-optimal feasible point. Corollary~\ref{face:endpoint} shows
that sparse endpoint probes can force exact termination. An explicit
$C^2$ example shows why the probes matter: ordinary backtracking can
continue taking half-steps forever.

A complementary analysis concerns the constrained
loss-dependent Hessian bound
\[
 \|\nabla^2 f(x)\|\le H_0+H_1\bigl(f(x)-f_Q^*\bigr),
 \qquad x\in Q.
\]
This condition and its function-dependent-smoothness context are
previously studied~\cite{warmup,liuwarmup,kale,caoselfbound,lobanov}. Under it,
Theorem~\ref{thm:fw} bounds the updates and all rejected function trials
needed by ordinary backtracking to reach its observable certificate. The
update bound separates the $H_0/\varepsilon$ and
$H_1\log(1/\varepsilon)$ contributions and includes the initial
estimate. Its lag potential controls curvature estimates that remain
large as the certificate falls; the method and full-region LMO are
unchanged. Theorem~\ref{lem:bridge} gives the sharp pointwise FW-gap
envelope, attained by one-dimensional objectives. For away-step FW,
Theorem~\ref{thm:afw} gives the corresponding convex bound, while
Theorem~\ref{thm:geometric} gives logarithmic certificate accuracy under
an additional geometric gap inequality. At $H_0=0$,
Proposition~\ref{prop:zeroface} shows that the full solution set is an
exposed face with a common gradient, and Theorem~\ref{thm:zerofinite}
proves finite exact termination of backtracking AFW without that
additional inequality. The finite identification mechanism uses
established AFW results~\cite{bomze,garber}; we do not claim a new AFW
algorithm or general identification theory.

A further certificate uses curvature along a segment rather than a
single worst-case bound. Propositions~\ref{vertex:radius} and
\ref{vertex:certificate} give an integrated-curvature gradient radius
and a verifiable optimal-vertex test with deterministic oracle errors.
The test requires a certified runner-up margin. It establishes
certificate validity, not a complexity guarantee for a noisy method.
Safe-screening methods for convex and generalized conditional-gradient
problems provide related certificate context~\cite{raj,sunbach}.

\paragraph{Related work.}
Harks~\cite[Theorem 1 and Corollary 3]{harks} proves last-iterate
distance and gap rates for FW on monotone variational inequalities,
assuming Jacobian coercivity along the exposed face at a unique solution
and using prescribed reciprocal-step schedules. Our face result instead
allows a full, possibly non-singleton minimizing face and assumes only a
Lipschitz gradient; it bounds a trajectory-wise gap-to-objective ratio
and the local certificate cost under ordinary backtracking. Harks's
open-loop rates do not give the endpoint-probe stopping result.

Vyguzov and Stonyakin~\cite{vyguzov} analyze FW with
$\|\nabla^2 f(x)\|\le L_0+L_1\|\nabla f(x)\|$ and an adaptive step
rule. Their curvature depends on gradient norm, whereas ours depends on
loss suboptimality. On compact $Q$, convexity gives
$f-f_Q^*\le D\|\nabla f\|$, but replacing this by a global maximum
gradient does not retain our separated dependence. Oikonomou and Loizou
derive a related whole-space self-bounding inequality by ray comparison
\cite[Lemma B.7, Eq.~(30)]{oikonomouloizou}; that auxiliary inequality
is not a contribution claimed here.

Backtracking trial counts, Fenchel dual certificates, and linear AFW
convergence are already established by Pedregosa et al.
\cite[Theorems 1, 3, and 4]{pedregosa}. Our scalar lower-bound maximum
is distinct from their averaged-gradient Fenchel dual. Computable
tangent bounds are also classical~\cite[Method 1]{freundgrigas};
optimistic and primal--dual short steps explicitly optimize gaps under
global smoothness~\cite[Proposition 4.1]{beyondshort}. Appendix
\ref{app:comparison} discusses these and other adaptive, screening, and
active-set comparisons. We make no minimax-optimality, universal
speedup, or absolute-priority claim. The experiments isolate
initialization, active-face geometry, and certificate cost; the
cache-controlled endpoint study found no accepted probes or practical
improvement on tasks whose optima are not vertices.
